\documentclass[11pt,a4paper]{article}
\usepackage[utf8]{inputenc}
\usepackage[T1]{fontenc}
\usepackage{geometry}
\geometry{margin=2.5cm}
\usepackage{graphicx}
\usepackage{hyperref}
\usepackage{listings}
\usepackage{xcolor}
\usepackage{enumitem}
\usepackage{fancyhdr}
\usepackage{lastpage}

\hypersetup{
    colorlinks=true,
    linkcolor=blue,
    urlcolor=blue,
}

% Solidity style
\lstdefinestyle{solidity}{
    language=C++,
    basicstyle=\footnotesize\ttfamily,
    keywordstyle=\color{blue},
    commentstyle=\color{gray},
    stringstyle=\color{green!50!black},
    numbers=left,
    numberstyle=\tiny,
    frame=single,
    breaklines=true,
    captionpos=b,
}

\pagestyle{fancy}
\fancyhf{}
\fancyhead[L]{ChainOfContractsLab}
\fancyhead[R]{Blockchain Systems}
\fancyfoot[C]{\thepage\ of \pageref{LastPage}}

\title{\textbf{Chain of Contracts Lab} \\\large Generic Multi-Contract Systems on Ethereum and Polkadot}
\author{Blockchain Systems Course}
\date{Spring 2026}

\begin{document}

\maketitle
\thispagestyle{empty}

\vspace{1cm}
\begin{center}
\fbox{\parbox{0.9\textwidth}{
\textbf{Learning Objectives:}
\begin{itemize}
    \item Understand the generic linked-list pattern for cross-contract communication
    \item Deploy and wire an arbitrarily long chain of smart contracts
    \item Verify chain execution through event analysis
    \item Compare contract models across Ethereum and Polkadot
\end{itemize}
}}
\end{center}
\vspace{1cm}

\newpage
\tableofcontents
\newpage

\section{Introduction}

In this lab, you will build a \textbf{generic chain of contracts} --- a
linked list of smart contracts where each contract calls the next, forming
a verifiable chain of \textbf{arbitrary length}. This pattern is fundamental
to many real-world blockchain applications: supply chain tracking, multi-step
approval workflows, layered DeFi protocols, and cross-chain message relays.

Unlike the fixed three-contract pattern (A $\to$ B $\to$ C) used in previous
versions, this lab uses a \textbf{generic linked-list architecture}:
\[
\text{User} \longrightarrow \textbf{ChainEntry} \longrightarrow \textbf{ChainLink(1)} \longrightarrow \textbf{ChainLink(2)} \longrightarrow \ldots \longrightarrow \textbf{ChainTerminal}
\]

You can insert any number of \texttt{ChainLink} contracts between the Entry
and Terminal, limited only by the blockchain's gas/weight limits.

You will implement this pattern on \textbf{two platforms}:
\begin{enumerate}
    \item \textbf{Ethereum} --- using Solidity and the Remix IDE
    \item \textbf{Polkadot} --- using ink! 5.0 and the Contracts UI
\end{enumerate}

Both implementations will be verified by an automated grader that checks event
ordering and final state correctness. The professor can also verify contracts
deployed on public testnets via Etherscan and Subscan.

\section{The Generic Chain Pattern}

\subsection{Architecture}

The pattern involves three contract roles that form a linked-list pipeline:

\[
\text{Entry} \longrightarrow \textbf{Link}_1 \longrightarrow \textbf{Link}_2 \longrightarrow \ldots \longrightarrow \textbf{Link}_N \longrightarrow \textbf{Terminal}
\]

\begin{itemize}
    \item \textbf{ChainEntry (Entry Point)}: Receives initial data from the user,
          records a receipt, emits \texttt{ChainStarted} and \texttt{EntryExecuted},
          and forwards to the first ChainLink via \texttt{processAndForward}.
    \item \textbf{ChainLink (Processor)}: Validates the data (non-empty,
          $\le$ 1024 bytes), prepends a \texttt{[VALIDATED]} marker, emits
          \texttt{LinkExecuted}, and forwards to the next contract via
          \texttt{receiveFinal}. Each link has a unique \texttt{stepNumber}.
    \item \textbf{ChainTerminal (Final Store)}: Accepts only from the last
          ChainLink, stores a \texttt{FinalRecord} with hash and timestamp,
          emits \texttt{FinalReceived} and \texttt{ChainCompleted}.
\end{itemize}

Each step emits events that the professor's grader will verify.
For a chain with one link (the default lab configuration), the event order is:
\begin{verbatim}
ChainStarted -> EntryExecuted -> LinkExecuted -> FinalReceived -> ChainCompleted
\end{verbatim}

For an N-link chain, \texttt{LinkExecuted} appears N times:
\begin{verbatim}
ChainStarted -> EntryExecuted -> LinkExecuted(1) -> ... -> LinkExecuted(N) -> FinalReceived -> ChainCompleted
\end{verbatim}

\subsection{Data Flow}

\begin{center}
\fbox{\parbox{0.8\textwidth}{
\textbf{Generic Chain Flow (1-link example)} \\
\vspace{0.3cm}

\noindent
1. User calls \texttt{ChainEntry.startChain(data)} \\
2. ChainEntry emits \texttt{ChainStarted}, \texttt{EntryExecuted}, calls
   \texttt{ChainLink.processAndForward(chainId, data)} \\
3. ChainLink validates, prepends \texttt{[VALIDATED]}, emits
   \texttt{LinkExecuted}, calls \texttt{ChainTerminal.receiveFinal(chainId, transformedData)} \\
4. ChainTerminal stores data, emits \texttt{FinalReceived} and
   \texttt{ChainCompleted} \\
5. Control returns to Entry, which marks the receipt as completed
}}
\end{center}

\subsection{Smart Contract Interfaces}

\paragraph{IChainEntry:}
\begin{lstlisting}[style=solidity]
interface IChainEntry {
    event ChainStarted(uint256 indexed chainId, address indexed initiator, bytes data);
    event EntryExecuted(uint256 indexed chainId, address indexed nextLink);

    function setNextLink(address _nextLink) external;
    function startChain(bytes calldata data) external returns (uint256);
    function getReceipt(uint256 _chainId) external view
        returns (uint256, address, bytes memory, bool);
    function getAllChainIds() external view returns (uint256[] memory);
}
\end{lstlisting}

\paragraph{IChainLink:}
\begin{lstlisting}[style=solidity]
interface IChainLink {
    event LinkExecuted(uint256 indexed chainId, uint256 indexed stepNumber,
        address indexed nextLink, bytes transformedData);
    event LinkValidationFailed(uint256 indexed chainId, string reason);

    function setNextLink(address _nextLink) external;
    function setPrevLink(address _prevLink) external;
    function setStepNumber(uint256 _step) external;
    function processAndForward(uint256 chainId, bytes calldata data) external returns (bool);
    function getStepNumber() external view returns (uint256);
    function getNextLink() external view returns (address);
    function getExecutionCount() external view returns (uint256);
}
\end{lstlisting}

\paragraph{IChainTerminal:}
\begin{lstlisting}[style=solidity]
interface IChainTerminal {
    event FinalReceived(address indexed from, uint256 indexed chainId,
        bytes data, uint256 timestamp);
    event ChainCompleted(uint256 indexed chainId, bytes32 dataHash);

    function setPrevLink(address _prevLink) external;
    function receiveFinal(uint256 chainId, bytes calldata data) external returns (bool);
    function getRecord(uint256 chainId) external view
        returns (uint256, bytes memory, bytes32, uint256, bool);
    function getRecordCount() external view returns (uint256);
}
\end{lstlisting}

\section{Ethereum Lab (Remix + Solidity)}

\subsection{Setup}

\begin{enumerate}
    \item Open \url{https://remix.ethereum.org/}
    \item Create a new workspace named \texttt{ChainOfContracts}
    \item Create the following files in the \texttt{contracts/} folder:
    \begin{itemize}
        \item \texttt{IChainEntry.sol}
        \item \texttt{IChainLink.sol}
        \item \texttt{IChainTerminal.sol}
        \item \texttt{ChainEntry.sol}
        \item \texttt{ChainLink.sol}
        \item \texttt{ChainTerminal.sol}
    \end{itemize}
\end{enumerate}

\subsection{Task 1: Implement ChainTerminal (Final Store)}

ChainTerminal is the simplest --- it stores validated data. Implement:

\begin{itemize}
    \item A \texttt{struct FinalRecord} containing: \texttt{chainId}, \texttt{data},
          \texttt{dataHash} (keccak256 of data), \texttt{timestamp}, \texttt{exists}.
    \item A \texttt{mapping(uint256 => FinalRecord)} for records.
    \item \texttt{setPrevLink(address)} --- only owner, one-time set.
    \item \texttt{receiveFinal(uint256, bytes)} --- only \texttt{prevLink} can call.
          Store the record, emit \texttt{FinalReceived} and \texttt{ChainCompleted}.
    \item \texttt{getRecord(uint256)} and \texttt{getRecordCount()} --- public queries.
\end{itemize}

\subsection{Task 2: Implement ChainLink (Processor)}

ChainLink validates incoming data and forwards to the next contract:

\begin{itemize}
    \item State: \texttt{nextLink}, \texttt{prevLink}, \texttt{owner},
          \texttt{stepNumber}, \texttt{executionCount}.
    \item Constructor takes \texttt{uint256 \_stepNumber} to identify this link's position.
    \item Access control: only \texttt{prevLink} can call \texttt{processAndForward}
          and \texttt{receiveFinal}.
    \item Validation rules:
    \begin{itemize}
        \item Data must not be empty.
        \item Data must not exceed 1024 bytes.
    \end{itemize}
    \item On success: prepend \texttt{[VALIDATED]} to data, emit
          \texttt{LinkExecuted}, call \texttt{receiveFinal} on the next contract.
    \item On failure: emit \texttt{LinkValidationFailed}, return false.
    \item Forwarding: use a low-level \texttt{call} with signature
          \texttt{receiveFinal(uint256,bytes)} so the same code works whether
          the next contract is another ChainLink or the ChainTerminal.
\end{itemize}

\subsection{Task 3: Implement ChainEntry (Entry Point)}

ChainEntry initiates the chain:

\begin{itemize}
    \item State: \texttt{nextLink}, \texttt{owner}, \texttt{chainCounter},
          \texttt{receipts} mapping.
    \item \texttt{startChain(bytes)}: increment counter, store receipt, emit
          \texttt{ChainStarted} and \texttt{EntryExecuted}, call
          \texttt{IChainLink(nextLink).processAndForward(...)}.
    \item Require \texttt{nextLink != address(0)} and non-empty data.
\end{itemize}

\subsection{Task 4: Deploy and Wire}

Follow the deployment order (reverse dependency):

\begin{enumerate}
    \item Deploy \texttt{ChainTerminal} $\to$ copy address T
    \item Deploy \texttt{ChainLink} with \texttt{stepNumber = 1}
          $\to$ call \texttt{setNextLink(T)} $\to$ copy address L
    \item Deploy \texttt{ChainEntry} $\to$ call \texttt{setNextLink(L)} $\to$ copy address E
    \item On \texttt{ChainLink}: call \texttt{setPrevLink(E)}
    \item On \texttt{ChainTerminal}: call \texttt{setPrevLink(L)}
\end{enumerate}

\subsection{Bonus: N-Link Chain}

For extra credit, add a second ChainLink:

\begin{enumerate}
    \item Deploy a second \texttt{ChainTerminal}
    \item Deploy \texttt{ChainLink(step=2)} $\to$ \texttt{setNextLink(Terminal)}
    \item Modify the existing \texttt{ChainLink(step=1)}: \texttt{setNextLink(Link2)}
    \item Wire: \texttt{Link2.setPrevLink(Link1)}, \texttt{Terminal.setPrevLink(Link2)}
    \item The event order becomes:
          \texttt{ChainStarted $\to$ EntryExecuted $\to$ LinkExecuted(1) $\to$ LinkExecuted(2) $\to$ FinalReceived $\to$ ChainCompleted}
    \item The stored data will be \texttt{[VALIDATED][VALIDATED]original}
\end{enumerate}

\subsection{Task 5: Execute and Verify}

\begin{enumerate}
    \item Call \texttt{ChainEntry.startChain("0x48656c6c6f")} (hex for "Hello")
    \item Check the transaction log for all 5 events in order
    \item Call \texttt{ChainTerminal.getRecord(1)} to verify the stored data
    \item The stored data should start with \texttt{[VALIDATED]}
\end{enumerate}

\section{Polkadot Lab (ink! 5.0)}

\subsection{Setup}

\begin{enumerate}
    \item Install Rust nightly:
    \begin{verbatim}
rustup install nightly
rustup default nightly
\end{verbatim}
    \item Install \texttt{cargo-contract}:
    \begin{verbatim}
cargo install cargo-contract --force
\end{verbatim}
    \item Install and run a local contracts node:
    \begin{verbatim}
cargo install contracts-node --git \
    https://github.com/paritytech/substrate-contracts-node.git
substrate-contracts-node --dev --tmp
\end{verbatim}
    \item Open \url{https://contracts-ui.substrate.io/} and connect to
          \texttt{ws://127.0.0.1:9944}
\end{enumerate}

\subsection{Task 6: Implement ChainTerminal (ink!)}

Create \texttt{polkadot/contracts/chain\_terminal/} with:

\begin{itemize}
    \item \texttt{\#[ink(storage)]} struct with: \texttt{prev\_link}, \texttt{owner},
          \texttt{record\_count}, \texttt{records: Mapping<u32, FinalRecord>}.
    \item \texttt{FinalRecord} struct with: \texttt{chain\_id}, \texttt{data},
          \texttt{data\_hash}, \texttt{timestamp}, \texttt{exists}.
    \item Events: \texttt{FinalReceived}, \texttt{ChainCompleted}.
    \item \texttt{receive\_final(chain\_id, data)} --- only \texttt{prev\_link} can call.
          Use \texttt{self.env().hash\_bytes::<ink::H256>(\&data)} for the hash.
    \item \texttt{set\_prev\_link(address)} --- only owner.
\end{itemize}

\subsection{Task 7: Implement ChainLink (ink!)}

Create \texttt{chain\_link/} with:

\begin{itemize}
    \item Constructor: \texttt{new(step\_number: u32)}.
    \item Access control: \texttt{assert\_eq!(caller, prev\_link)}.
    \item Both \texttt{process\_and\_forward} and \texttt{receive\_final} serve
          as entry points (so the previous link can call either).
    \item Cross-contract forward using \texttt{build\_call}:
\begin{lstlisting}[style=rust]
build_call::<DefaultEnvironment>()
    .call(next_link)
    .exec_input(
        ExecutionInput::new(
            Selector::new(ink::selector_bytes!("receive_final"))
        )
        .push_arg(chain_id)
        .push_arg(transformed_data),
    )
    .returns::<bool>()
    .invoke()
\end{lstlisting}
\end{itemize}

\subsection{Task 8: Implement ChainEntry (ink!)}

Create \texttt{chain\_entry/} with:

\begin{itemize}
    \item Counter-based \texttt{chain\_counter} incremented on each call.
    \item Store \texttt{ChainReceipt} before forwarding.
    \item Cross-contract call to Link using \texttt{build\_call}
          with the \texttt{process\_and\_forward} selector.
    \item Update receipt to \texttt{completed = true} after successful call.
\end{itemize}

\subsection{Task 9: Deploy and Wire (Contracts UI)}

Same order as Ethereum:

\begin{enumerate}
    \item Build each contract: \texttt{cargo contract build --release}
    \item Upload \& instantiate ChainTerminal
    \item Upload \& instantiate ChainLink(step=1)
          $\to$ call \texttt{setNextLink(terminalAddr)}
    \item Upload \& instantiate ChainEntry
          $\to$ call \texttt{setNextLink(linkAddr)}
    \item On ChainLink: call \texttt{setPrevLink(entryAddr)}
    \item On ChainTerminal: call \texttt{setPrevLink(linkAddr)}
    \item Call \texttt{ChainEntry.startChain("0x48656c6c6f")}
\end{enumerate}

\section{Submission Requirements}

Your submission must include:

\subsection{Ethereum (50\% of Grade)}
\begin{itemize}
    \item \texttt{ChainEntry.sol}, \texttt{ChainLink.sol}, \texttt{ChainTerminal.sol}
    \item All three files in a folder named \texttt{YourName-eth/}
\end{itemize}

\subsection{Polkadot (50\% of Grade)}
\begin{itemize}
    \item \texttt{chain\_entry/lib.rs}, \texttt{chain\_link/lib.rs}, \texttt{chain\_terminal/lib.rs}
    \item All three files in a folder named \texttt{YourName-polkadot/}
\end{itemize}

\subsection{Grading Rubric}

\begin{itemize}
    \item \textbf{Compilation (20\%)}: Contracts compile without errors.
    \item \textbf{Deployment (30\%)}: All three contracts deploy and wire correctly.
    \item \textbf{Event Chain (30\%)}: Events emitted in correct order:
          \texttt{ChainStarted $\to$ EntryExecuted $\to$ LinkExecuted $\to$ FinalReceived $\to$ ChainCompleted}.
    \item \textbf{Final State (20\%)}: ChainTerminal stores the record with correct
          \texttt{[VALIDATED]} prefix.
\end{itemize}

Passing threshold: 70\%

\section{Troubleshooting}

\begin{itemize}
    \item \textbf{Compilation errors}: Check Solidity version (0.8.24),
          import paths, and interface function signatures.
    \item \textbf{"NextLink not set"}: Did you call \texttt{setNextLink} on ChainEntry?
    \item \textbf{"Only prevLink can call this"}: Did you call \texttt{setPrevLink}
          on the ChainLink (set to Entry's address)?
    \item \textbf{"Only prevLink can call this" on Terminal}: Did you call
          \texttt{setPrevLink} on ChainTerminal (set to Link's address)?
    \item \textbf{Missing events}: Check that each function calls \texttt{emit EventName}.
    \item \textbf{ink! build fails}: Ensure you're using Rust nightly and
          the latest \texttt{cargo-contract}.
    \item \textbf{Cross-contract call fails}: Verify the downstream contract's
          address and function selector. The forwarding uses \texttt{receiveFinal}
          as the uniform selector.
    \item \textbf{Chain forwarding failed}: The nextLink address may be wrong
          or the downstream contract's \texttt{receiveFinal} reverted.
\end{itemize}

\section{Reference Materials}

\begin{itemize}
    \item Solidity Documentation: \url{https://docs.soliditylang.org/}
    \item Remix IDE: \url{https://remix.ethereum.org/}
    \item ink! Documentation: \url{https://use.ink/}
    \item Contracts UI: \url{https://contracts-ui.substrate.io/}
    \item TON Comparison (provided separately)
    \item Etherscan API: \url{https://docs.etherscan.io/api-endpoints/logs}
    \item Subscan API: \url{https://docs.api.subscan.io/}
\end{itemize}

\vspace{1cm}
\begin{center}
\fbox{\parbox{0.9\textwidth}{
\textbf{Good luck!} Remember that the chain must produce all 5 events in
exact order for the grader to give full credit. Test your chain by calling
\texttt{startChain} and checking the full transaction receipt before submitting.
For bonus credit, try adding a second ChainLink to test the N-link pattern.
}}
\end{center}

\end{document}